\section{Introduction}
\label{sec:introduction}
\IEEEPARstart{P}{hysical} intelligence---agents whose perception, decision, and action loops close in the physical world under tight energy and latency budgets---demands controllers that remain robust on imperfect hardware\cite{yuan2026embodiedr,guo2026neurovla}. Neuromorphic systems promise the required efficiency\cite{katti2026neuromorphic,gerven2025neuromorphic}, and spiking implementations on Loihi~2 and FPGA substrates now cover vision and control workloads\cite{wang2025neuromorphic,arfa2025hardwareaware,gautam2025neurocorex}. What remains poorly quantified is how \emph{wiring} determines robustness: engineered spiking networks are designed task-by-task, so their fault tolerance must be re-measured per design\cite{thieu2022effects,yunusoglu2025neuromorphic}. Complete, signed wiring diagrams of whole brains are now public, and point-neuron models built on them reproduce realistic dynamics without training\cite{shiu2024drosophila,lin2024network}; an evolved connectome is thus a field-tested controller and a reference point for engineered fault tolerance.

The released leaky integrate-and-fire (LIF) whole-brain model of \emph{Drosophila melanogaster}\cite{shiu2024drosophila}\footnote{\url{https://github.com/philshiu/Drosophila_brain_model}}, which instantiates the complete FlyWire connectome\cite{lin2024network}, already used targeted computational silencings to identify circuit neurons\cite{shiu2024drosophila}, and we extend that perturbation style into a systematic sweep across \emph{hardware-motivated fault classes}, scored jointly for transmission and activity-envelope compliance. We instrument the visuomotor escape pathway: looming-sensitive visual interneurons (LC4, LPLC2) converge onto descending command neurons (DNp01, DNa02, DNp04, DNb01) whose activation evokes take-off and escape steering\cite{simpson2023descending}. Its failure modes map directly onto hardware fault classes: weight imprecision, broken interconnect, dead cores, dead sensors.

We inject five fault families (Section~III-C), each over 12 seeded 600\,ms trials, with two controls: an edge-permuted null preserving weights and degree statistics, and an analytic pathway-mass computation. Our contributions: (i) a fault-injection framework mapping hardware failure modes onto a connectome-constrained 138k-neuron LIF model; and (ii) the first tolerance curves of descending-neuron output in a complete connectome (50\% weight noise absorbed, 50\% connection loss retains 76\%, 10\% unit loss leaves the median unchanged): top-1\% hub loss instead produces a $2.5\times$ \emph{activity runaway} that passes transmission while violating the output envelope, and whole-channel deafferentation is the only catastrophic mode. We further (iii) decompose why graph redundancy alone explains neither specificity nor graceful degradation (threshold amplification as a candidate account) and (iv) distill five candidate design implications for fault-tolerant controllers.
